\section{Requirements}
\label{sec:requirements}
%% ====================================================================

The ontology requirements were gathered through iterative discussions with
domain experts (materials scientists and MLIP developers) across the seven work
packages of the META-LEARN project~\cite{grabowski2024metalearnerc}. We
followed an iterative ontology engineering methodology, identifying the
following competency questions that the ontology should be able to answer:

\begin{enumerate}
\item[CQ1.] Which MLIP algorithms exist, and what hyperparameters does each
  accept (with types, ranges, and default values)?
\item[CQ2.] Which libraries implement a given algorithm, and what versions
  are available?
\item[CQ3.] For a given material system, which training datasets are available,
  and what DFT code and settings were used to generate them?
\item[CQ4.] What is the provenance of a training dataset (published source,
  in-house calculation, augmented from existing data)?
\item[CQ5.] How many atomic configurations does a training dataset contain,
  and what properties are covered (energies, forces, stresses)?
\item[CQ6.] Which published studies have benchmarked a given algorithm on a
  given material property, and what accuracy metrics were reported?
\item[CQ7.] For a given material and target property, how do different
  algorithm and hyperparameter combinations rank by accuracy?
\item[CQ8.] Which simulation types (molecular dynamics, Monte Carlo, geometry
  optimization, phonon calculations, thermodynamic integration) have been
  performed with a given MLIP?
\item[CQ9.] For a given target accuracy, which algorithms and trained
  models are most efficient---both in terms of asymptotic complexity
  (prior knowledge) and measured training and inference cost
  (empirical data)?
\end{enumerate}

Additionally, we identified the following design requirements:

\begin{enumerate}
\item[DR1.] \emph{Alignment with existing ontologies.} Reuse classes and
  properties from MDO, CMSO/ASMO, EMMO, ML-Schema, and PROV-O where
  appropriate, to maximize interoperability.
\item[DR2.] \emph{SHACL validation.} Define SHACL shapes for each entity
  type to enforce data quality and enable form-based data entry.
\item[DR3.] \emph{Linked Data integration.} Link entities to Wikidata where
  possible (chemical elements, materials, algorithms, researchers).
\item[DR4.] \emph{Extensibility.} The ontology should be extensible to
  accommodate new algorithm families and hyperparameters as the field evolves.
\end{enumerate}

%% ====================================================================
